\documentclass[10pt,a4paper]{article}
% MS2 summary · compile with pdflatex (Overleaf default) · output: docs/ms2_summary.pdf

\usepackage[a4paper,margin=1.7cm,top=1.5cm,bottom=1.6cm]{geometry}
\usepackage[T1]{fontenc}
\usepackage[utf8]{inputenc}
\usepackage[scaled=0.95]{helvet}
\renewcommand{\familydefault}{\sfdefault}
\usepackage{microtype}
\usepackage[table]{xcolor}
\usepackage{tabularx,booktabs,array}
\usepackage{enumitem}
\usepackage{titlesec}
\usepackage[hidelinks]{hyperref}

\definecolor{accent}{HTML}{1F4E79}
\definecolor{changed}{HTML}{1E7B34}
\definecolor{partly}{HTML}{B26A00}
\definecolor{rowgray}{HTML}{F3F5F7}

\titleformat{\section}{\large\bfseries\color{accent}}{\thesection}{0.6em}{}
\titlespacing*{\section}{0pt}{1.0em}{0.4em}
\titleformat{\subsection}{\normalsize\bfseries}{}{0em}{}
\titlespacing*{\subsection}{0pt}{0.7em}{0.25em}
\setlength{\parindent}{0pt}
\setlength{\parskip}{0.35em}
\setlist{nosep,leftmargin=1.2em,topsep=0.2em}
\renewcommand{\arraystretch}{1.25}
\newcolumntype{Y}{>{\raggedright\arraybackslash}X}
\newcolumntype{L}[1]{>{\raggedright\arraybackslash}p{#1}}
\newcommand{\file}[1]{\texttt{\small #1}}
\newcommand{\changed}{\textbf{\color{changed}Changed}}
\newcommand{\partly}{\textbf{\color{partly}Partly}}

\begin{document}

{\LARGE\bfseries\color{accent} MS2 Summary · Feature Pipeline}\\[0.25em]
{\large P04 · Swiss Segmented Unemployment Rate Predictor}\\[0.15em]
{\small MLOps HS26 · Repository: \url{https://github.com/theodorahaimoff/HSLU_HS26_MLOPS} · Decision log: \file{docs/decisions.md} (D1--D13)}

\section{What I achieved in this milestone}

The feature pipeline runs on a schedule without manual steps. Every day it pulls three live sources, stores every
download with its date and hash, rebuilds a point-in-time correct feature table for all 45 series (national total
and 44 gender~×~nationality~×~age segments) and writes it as a new version to a versioned Parquet feature store on GCS.

\rowcolors{2}{rowgray}{white}
\begin{tabularx}{\textwidth}{@{}L{3.6cm}Y@{}}
\toprule
\textbf{Requirement} & \textbf{What works now} \\
\midrule
Live ingestion + backfill &
SAKE (BFS, quarterly target) via the opendata.swiss CKAN API, SECO Weekly Economic Activity (weekly) and LIK
inflation from the SNB data portal (monthly). Each run recomputes the full backfill: one row set per SAKE release
since 2011Q1 (62~release dates~× 45~series, the MS3 training data) plus \emph{refresh} rows as of the run date, the
current nowcast. \\
Features without leakage &
Every value gets an \texttt{available\_from} date (release date, or period end + publication lag). A row may only
use values available on its \texttt{as\_of} date. Three checks run before anything is written; the strongest
rebuilds random dates after deleting all later data and requires identical features. \\
Feature store &
GCS bucket: \file{raw/pull\_log.csv} (download date, SHA-256, release date), raw files stored only when their
content changed, and \file{feature\_store/features/version=<date>/} with \file{features.parquet} and a
\file{manifest.json} (source hashes, git commit, lags). \\
Scheduled execution &
GitHub Actions \file{feature\_pipeline.yml}, daily 05:17 UTC (weekly until 02.11, D12). Scheduled runs on 26.10 and
02.11, daily from 03.11, all green; tests run first, so a broken pipeline writes no version. \\
Docker, pinned deps, README &
\file{Dockerfile}: \texttt{python:3.14-slim}, uv pinned, dependency layer first, non-root user, credentials only at
run time. All installs use \texttt{uv sync -{}-locked} on \file{uv.lock}. README: architecture, data, feature store
layout, how to run locally, in Docker and in CI. \\
Tests, CI, secrets &
25 unit tests on synthetic files (no network, no credentials); \file{ci.yml} runs lint, tests and the Docker build on
every push. Keys live in \file{.env} (gitignored) and GitHub Secrets; the history contains none. \\
\bottomrule
\end{tabularx}

\subsection{Key decisions}
\begin{itemize}
  \item \textbf{Nowcast instead of a static quarterly pipeline} (D1--D3): SAKE stays the target, weekly WEA and monthly
        inflation make the features change between releases.
  \item \textbf{Target = year-on-year change} (D6): SAKE is not seasonally adjusted (Q1 and Q3 lie about 0.3~pp above the yearly mean, Q2 and Q4 below).
  \item \textbf{Point-in-time features at the SAKE release date} (D7), with publication lags observed in the data
        (SAKE 49~days, WEA $\le$24, LIK 21) and chosen on the late side: too late is pessimistic, too early would leak.
  \item \textbf{Versioning in the pipeline instead of DVC} (D11): data arrives on a schedule, so every pull is hashed
        and stored at download time, and every feature version names its source hashes and git commit.
\end{itemize}
Observed on the way: SECO revised past WEA values between my pulls of 15.10 and 25.10 (2026Q1 mean $-0.41 \rightarrow
-0.49$), which is exactly why every pull is kept with its date.

\subsection{What a reviewer should look at first}
\begin{enumerate}
  \item \file{notebooks/02\_feature\_pipeline.ipynb}: every step of the pipeline with its outputs and findings.
  \item \file{src/feature\_pipeline/features.py}: \texttt{build\_features}, \texttt{backfill} and the leakage checks
        in \texttt{check\_features}; proven by \file{tests/test\_features.py::test\_no\_leakage\_on\_any\_date}.
  \item The \textbf{Actions} tab: scheduled runs of \emph{Feature pipeline}, and \emph{ci} on every push.
  \item Run it yourself, no credentials needed: \texttt{uv sync -{}-locked \&\& uv run python -m
        src.feature\_pipeline.run} pulls the live sources and writes the feature store to \file{data/}.
\end{enumerate}

\newpage
\section{Response to reviews}
Every labelled point from both MS1 reviews and the editorial summary. Details and alternatives: \file{docs/decisions.md}.

\rowcolors{2}{rowgray}{white}
\begin{tabularx}{\textwidth}{@{}L{0.9cm}L{3.1cm}L{1.35cm}Y@{}}
\toprule
\textbf{Label} & \textbf{Point} & \textbf{Status} & \textbf{Answer · where} \\
\midrule
E1 \newline \small(Must) & Data too static: SAKE changes once a quarter, so ingest $\rightarrow$ retrain $\rightarrow$ monitor
can hardly be shown & \changed &
Nowcast design: SAKE stays the quarterly target; SECO WEA (weekly) and monthly LIK inflation were added as live
features. New WEA or LIK values reach the feature store as refresh rows on the next daily run, so the forecast input
changes every week. A monthly SECO target was rejected: no API, monthly figures only as PDF (would need scraping).
\newline\file{docs/decisions.md} D1--D3, D12 · \file{src/feature\_pipeline/sources.py} \\
E2 & Weekly scheduled retraining on unchanged quarterly labels produces identical models & \partly &
Decided: retrain on a new SAKE release, not on a schedule (D4). The trigger signal exists now: a new SAKE hash in the
pull log sets \texttt{new\_sake\_release} in the manifest of that feature version. The retraining workflow that reacts
to it is MS3 work. \newline\file{src/feature\_pipeline/pull\_log.py} \\
A1 & Release detection by the hash of one fixed file URL misses releases published as new resources; store release
dates & \changed &
Each run lists the dataset's resources through CKAN and takes the CSV with the latest \texttt{issued} date, not a fixed
URL. The pull log stores download date, release date, latest quarter and SHA-256 per pull. Tested: a newer release
as a separate resource is picked (\file{tests/test\_sources.py}), and two consecutive pulls store only what changed
(\file{tests/test\_storage.py}). The first real new release (2026Q3, expected around 18.11) is the live check. \newline D5 · \file{src/feature\_pipeline/sources.py} (\texttt{pull\_sake}) \\
A2, B1 & Small segments: how does BFS mark unreliable or suppressed values, and how is it shown? & \partly &
BFS marks reliability in a separate \texttt{STATUS} column (A normal, U low reliability, O missing), not with brackets.
In scope there are no O values, but 17 of 45 segments have U quarters (7 in more than half). The code is kept as
feature \texttt{status\_t}; the success target is measured on the 28 segments without U and the others are reported
separately. The UI flag follows with the UI in MS4. \newline D9 · \file{notebooks/01\_data\_exploration.ipynb} §2--3 \\
B2 & 15\,\% target against seasonal naive only; persistence may be the stronger baseline & \changed &
Both computed on the stored features, last 12 target quarters, MAE in pp: seasonal naive 0.794 vs persistence 0.976
over all 45 series, and 0.630 vs 0.672 on the 28 reliable segments. Seasonal naive is stronger in every group
(persistence carries last quarter's season into the next). Restated target: MAE $\le$ 0.535~pp on the 28 reliable
segments, 15\,\% below seasonal naive. \newline D10 · \file{src/training\_pipeline/baselines.py} \\
B3 & \texttt{requirements.txt} with minimum versions only & \changed &
\file{uv.lock} pins every package with version and hash; CI, the scheduled job and the Docker build install with
\texttt{uv sync -{}-locked}, which fails on a stale lockfile. \file{requirements.txt} removed; dev tools in a separate
group, not in the image. \newline D8 · \file{pyproject.toml}, \file{uv.lock}, \file{Dockerfile} \\
E3 & Why FastAPI \emph{and} Streamlit? & \changed &
FastAPI dropped. The forecasts are batch-shaped: all inputs are in the feature store before anyone asks, and all
45 series are forecast per run, so an API would only return what the batch job already computed. A single Streamlit
app on Cloud Run will read the forecasts from GCS. Model version, feature time and creation time are stored per
forecast row and shown in the UI. \newline D13 · \file{docs/decisions.md} \\
\bottomrule
\end{tabularx}

\subsection{Next (MS3)}
Training pipeline reading the feature store: pooled LightGBM on the year-on-year target, walk-forward backtest over
the last 12 quarters against the seasonal-naive baseline, W\&B tracking and registry with a \emph{champion} alias, and
a retraining workflow triggered by a new SAKE release (E2).

\vfill
{\footnotesize\color{gray} AI use: I used Claude (Anthropic) as a coding and writing assistant for this milestone. I
reviewed, ran and can explain every part of the code and this summary.}

\end{document}
